\section{Gradients of Linear Layers in Transformers}
\label{sec:gradients_linear_layers}

We begin by formalizing the relationship between model parameters, activations, and
gradients for a linear projection in a transformer layer.
This relationship is central to understanding how low-rank and spiky structure
emerges in gradient matrices and how it influences optimization dynamics.

\paragraph{Linear projection and model parameters.}
Consider a single transformer layer and one of its linear projections, for example
the key projection matrix
\[
W_k \in \mathbb{R}^{d_{\text{out}} \times d_{\text{in}}}.
\]
In standard transformer architectures, $d_{\text{out}} = d_{\text{in}} = d$.
Given an input hidden state
\[
x \in \mathbb{R}^{d_{\text{in}}},
\]
the key vector is computed as
\[
k = W_k x \in \mathbb{R}^{d_{\text{out}}}.
\]

The rows of $W_k$ correspond to output features of the key vector, while the columns
correspond to input features of the hidden state.
The parameter matrix $W_k$ therefore defines how input features are linearly combined
to produce output features.

\paragraph{Token-level activations.}
During training, the model processes a batch of tokens.
For token $j$, let
\[
x^{(j)} \in \mathbb{R}^{d_{\text{in}}}
\]
denote the input activation to the key projection.
Stacking activations across a batch of $n$ tokens yields the activation matrix
\[
X \in \mathbb{R}^{n \times d_{\text{in}}},
\qquad
X_{j,:} = \bigl(x^{(j)}\bigr)^\top.
\]
The activation matrix encodes the representation geometry of the model at the input
of the linear layer.

\paragraph{Backpropagated error signals.}
Let
\[
\delta_k^{(j)} = \frac{\partial \mathcal{L}}{\partial k^{(j)}}
\in \mathbb{R}^{d_{\text{out}}}
\]
denote the backpropagated error signal associated with token $j$ at the key
projection.
Stacking these error vectors gives
\[
\Delta_k \in \mathbb{R}^{n \times d_{\text{out}}},
\qquad
(\Delta_k)_{j,:} = \bigl(\delta_k^{(j)}\bigr)^\top.
\]

\paragraph{Gradient matrix.}
During backpropagation, the gradient of the loss $\mathcal{L}$ with respect to $W_k$
has the same shape as the parameter matrix:
\[
G \;\triangleq\; \nabla_{W_k} \mathcal{L}
\;\in\;
\mathbb{R}^{d_{\text{out}} \times d_{\text{in}}}.
\]
For a single token $j$, the contribution to the gradient has the outer-product form
\[
\nabla_{W_k} \mathcal{L}^{(j)}
=
\delta_k^{(j)} \bigl(x^{(j)}\bigr)^\top
\;\in\;
\mathbb{R}^{d_{\text{out}} \times d_{\text{in}}}.
\]
Summing over all $n$ tokens in the batch, the full gradient matrix can be written as
\[
G
=
\sum_{j=1}^{n}
\delta_k^{(j)} \bigl(x^{(j)}\bigr)^\top
=
\Delta_k^\top X.
\]
Thus, although $G$ is a single $d_{\text{out}} \times d_{\text{in}}$ matrix, it
implicitly aggregates correlations between activation directions and error
directions across many tokens and contexts.

\paragraph{Connection to stochastic gradient descent.}
Under vanilla stochastic gradient descent (SGD), the parameter update at training
step $t$ is
\[
W_k^{(t+1)}
=
W_k^{(t)} - \eta \, G^{(t)},
\]
where $\eta > 0$ is the learning rate.
Substituting the gradient expression yields
\[
W_k^{(t+1)}
=
W_k^{(t)}
-
\eta
\sum_{j=1}^{n}
\delta_k^{(j,t)} \bigl(x^{(j,t)}\bigr)^\top.
\]

Unrolling SGD over multiple steps shows that the parameter matrix can be interpreted
as a time-integrated sum of gradient outer products:
\[
W_k^{(T)}
=
W_k^{(0)}
-
\eta
\sum_{t=1}^{T}
\sum_{j}
\delta_k^{(j,t)} \bigl(x^{(j,t)}\bigr)^\top.
\]
Consequently, the structure of the parameter matrix reflects the accumulated
correlations between activation patterns and backpropagated error signals over the
course of training.

\paragraph{Summary.}
Activations $X$ describe what features are present in the data, gradients $G$ describe
how the model parameters are being pushed to change in response to those features,
and the parameter matrix $W_k$ encodes the historical integration of these gradient
signals.
In the following sections, we focus on the gradient matrix $G$ as the primary object
for analyzing optimization geometry and for designing simple, GPU-friendly
approximations to its dominant low-dimensional structure.